\documentclass[10pt,a4paper]{amsart}
\usepackage{amsmath,amssymb,amsthm,graphicx}
\usepackage{verbatim}
\newenvironment{reference}{\comment}{\endcomment}
%\newenvironment{reference}{}{}
\newenvironment{slogan}{\comment}{\endcomment}
\newenvironment{history}{\comment}{\endcomment}

% For commutative diagrams we use Xy-pic
\usepackage[all]{xy}

% We use 2cell for 2-commutative diagrams.
\xyoption{2cell}
% Theorem environments.
%
\theoremstyle{plain}
\newtheorem{theorem}[subsection]{Theorem}
\newtheorem{proposition}[subsection]{Proposition}
\newtheorem{lemma}[subsection]{Lemma}

\theoremstyle{definition}
\newtheorem{definition}[subsection]{Definition}
\newtheorem{example}[subsection]{Example}
\newtheorem{exercise}[subsection]{Exercise}
\newtheorem{situation}[subsection]{Situation}

\theoremstyle{remark}
\newtheorem{remark}[subsection]{Remark}
\newtheorem{remarks}[subsection]{Remarks}

\numberwithin{equation}{subsection}

% Macros
%
\def\lim{\mathop{\mathrm{lim}}\nolimits}
\def\colim{\mathop{\mathrm{colim}}\nolimits}
\def\Spec{\mathop{\mathrm{Spec}}}
\def\Hom{\mathop{\mathrm{Hom}}\nolimits}
\def\Ext{\mathop{\mathrm{Ext}}\nolimits}
\def\SheafHom{\mathop{\mathcal{H}\!\mathit{om}}\nolimits}
\def\SheafExt{\mathop{\mathcal{E}\!\mathit{xt}}\nolimits}
\def\Sch{\mathit{Sch}}
\def\Mor{\mathop{\mathrm{Mor}}\nolimits}
\def\Ob{\mathop{\mathrm{Ob}}\nolimits}
\def\Sh{\mathop{\mathit{Sh}}\nolimits}
\def\NL{\mathop{N\!L}\nolimits}
\def\CH{\mathop{\mathrm{CH}}\nolimits}
\def\proetale{{pro\text{-}\acute{e}tale}}
\def\etale{{\acute{e}tale}}
\def\QCoh{\mathit{QCoh}}
\def\Ker{\mathop{\mathrm{Ker}}}
\def\Im{\mathop{\mathrm{Im}}}
\def\Coker{\mathop{\mathrm{Coker}}}
\def\Coim{\mathop{\mathrm{Coim}}}

% Boxtimes
%
\DeclareMathSymbol{\boxtimes}{\mathbin}{AMSa}{"02}

%
% Macros for moduli stacks/spaces
%
\def\QCohstack{\mathcal{QC}\!\mathit{oh}}
\def\Cohstack{\mathcal{C}\!\mathit{oh}}
\def\Spacesstack{\mathcal{S}\!\mathit{paces}}
\def\Quotfunctor{\mathrm{Quot}}
\def\Hilbfunctor{\mathrm{Hilb}}
\def\Curvesstack{\mathcal{C}\!\mathit{urves}}
\def\Polarizedstack{\mathcal{P}\!\mathit{olarized}}
\def\Complexesstack{\mathcal{C}\!\mathit{omplexes}}
% \Pic is the operator that assigns to X its picard group, usage \Pic(X)
% \Picardstack_{X/B} denotes the Picard stack of X over B
% \Picardfunctor_{X/B} denotes the Picard functor of X over B
\def\Pic{\mathop{\mathrm{Pic}}\nolimits}
\def\Picardstack{\mathcal{P}\!\mathit{ic}}
\def\Picardfunctor{\mathrm{Pic}}
\def\Deformationcategory{\mathcal{D}\!\mathit{ef}}
\usepackage{cleveref}
\crefname{originallemma}{original lemma}{original lemmas}
\Crefname{originallemma}{Original lemma}{Original lemmas}
\makeatletter
\newlabel{algebra-lemma-pullback-module}{{02M1}{source}{}{sourcepointer-1}{}}
\newlabel{algebra-lemma-pullback-module@cref}{{[originallemma][0][]02M1}{[1][0][]source}{}{}{}}
\makeatother
\crefname{originallemma}{original lemma}{original lemmas}
\Crefname{originallemma}{Original lemma}{Original lemmas}
\makeatletter
\newlabel{decent-spaces-lemma-describe-henselian-local-ring}{{0EMY}{source}{}{sourcepointer-2}{}}
\newlabel{decent-spaces-lemma-describe-henselian-local-ring@cref}{{[originallemma][0][]0EMY}{[1][0][]source}{}{}{}}
\makeatother
\crefname{originalreference}{original reference}{original references}
\Crefname{originalreference}{Original reference}{Original references}
\makeatletter
\newlabel{definition-order-vanishing}{{02RJ}{source}{}{sourcepointer-3}{}}
\newlabel{definition-order-vanishing@cref}{{[originalreference][0][]02RJ}{[1][0][]source}{}{}{}}
\makeatother
\crefname{originaldefinition}{original definition}{original definitions}
\Crefname{originaldefinition}{Original definition}{Original definitions}
\makeatletter
\newlabel{divisors-definition-order-vanishing}{{02RJ}{source}{}{sourcepointer-4}{}}
\newlabel{divisors-definition-order-vanishing@cref}{{[originaldefinition][0][]02RJ}{[1][0][]source}{}{}{}}
\makeatother
\title{The order of vanishing at a prime divisor meeting the schematic locus agrees with the scheme definition}
\author{The Stacks Project Authors}
\date{stacks-project:spaces-over-fields.tex:local-snapshot-20261012; GNU FDL 1.2 or later}
\begin{document}
\maketitle
\noindent\textit{Editorial note.} All delivered mathematical text is the author's own, copied from the pinned Stacks Project chapter file spaces-over-fields.tex; no statement or proof step was rewritten, reordered or removed.\par\medskip
\noindent\textit{Editorial note.} The counted proof is the author's complete argument: the reduction to an integral Noetherian scheme, the identification of the henselian local ring as a henselization together with the restatement as an equality of lengths, and the conclusion from the length formula for a flat local extension.\par\medskip
\noindent\textit{Editorial note.} The label lemma-order-vanishing-agrees names the counted statement, and the proof environment immediately following it in the source is delivered verbatim as its proof.\par\medskip
\noindent\textit{Editorial note.} The counted statement is carried by the author's own lemma environment in the chapter source.\par\medskip
\noindent\textit{Editorial note.} Pre-open gate: python operations/agent-authoring/tools/precard\_gate.py --source spaces-over-fields.tex --label lemma-order-vanishing-agrees --start 876 --end 902 --local-id 28066 returned RESULT=OK (L28066 free, occupancy FREE, public interval clear). Candidate taken from sources/stacks-spaces-over-fields/candidates-halfA.json.\par\medskip
\noindent\textit{Editorial note.} The statement and proof are copied from the pinned chapter file; the cited inputs are kept as external pointers with their canonical Stacks tags.\par\medskip
\noindent\textit{Allowed standard prerequisites (editorial).}\begin{itemize}
\item Description of the henselian local ring at a point of the schematic locus. Decent Spaces, the lemma identifying the henselian local ring at a point lying in the schematic locus with the henselization of the local ring of that scheme. Supplies the identification of the henselian local ring used in the definition with a henselization of a Noetherian local ring.
\item Invariance of length under a flat local extension. Algebra, the lemma computing the length of a quotient after a flat base change by a local ring homomorphism, and the fact that a local ring maps flatly into its henselization. Supplies the equality of the two lengths appearing in the statement.
\item Covering of the schematic locus by affine etale neighbourhoods. Same chapter, the covering lemma for normal or schematic pieces used to reduce to the case of an integral Noetherian scheme. Supplies the reduction to the case where the space is an integral Noetherian scheme.
\end{itemize}\medskip
\ifdefined\hypertarget\hypertarget{sourcepointer-1}{}\fi
\noindent\textit{Source reference.} Original lemma 02M1: sources/stacks-chow/tags.txt, source lines 3428--3428. Prerequisite pointer invoked at proof line(s) 899; it names a standard fact of a sibling chapter that the counted argument uses, and every step of the counted proof built on it is retained verbatim in the delivered body.\par\medskip
\ifdefined\hypertarget\hypertarget{sourcepointer-2}{}\fi
\noindent\textit{Source reference.} Original lemma 0EMY: sources/stacks-chow/tags.txt, source lines 18323--18323. Prerequisite pointer invoked at proof line(s) 890; it names a standard fact of a sibling chapter that the counted argument uses, and every step of the counted proof built on it is retained verbatim in the delivered body.\par\medskip
\ifdefined\hypertarget\hypertarget{sourcepointer-3}{}\fi
\noindent\textit{Source reference.} Original reference 02RJ: sources/stacks-spaces-over-fields/pinned-original-source.tex, source lines 838--838. Prerequisite pointer invoked at proof line(s) 881; it names a standard fact of the same chapter that the counted argument uses, and every step of the counted proof built on it is retained verbatim in the delivered body.\par\medskip
\ifdefined\hypertarget\hypertarget{sourcepointer-4}{}\fi
\noindent\textit{Source reference.} Original definition 02RJ: sources/stacks-chow/tags.txt, source lines 3626--3626. Prerequisite pointer invoked at proof line(s) 883; it names a standard fact of a sibling chapter that the counted argument uses, and every step of the counted proof built on it is retained verbatim in the delivered body.\par\medskip
\begin{lemma}
\label{lemma-order-vanishing-agrees}
Let $S$ be a scheme. Let $X$ be a locally Noetherian integral algebraic
space over $S$. Let $f \in R(X)^*$. If the prime divisor
$Z \subset X$ meets the schematic locus of $X$, then the order
of vanishing $\text{ord}_Z(f)$ of Definition \ref{definition-order-vanishing}
agrees with the order of vanishing of
Divisors, Definition \ref{divisors-definition-order-vanishing}.
\end{lemma}

\begin{proof}
After shrinking $X$ we may assume $X$ is an integral Noetherian scheme.
If $\xi \in Z$ denotes the generic point, then we find that
$\mathcal{O}_{X, \xi}^h$ is the henselization of $\mathcal{O}_{X, \xi}$
(Decent Spaces, Lemma \ref{decent-spaces-lemma-describe-henselian-local-ring}).
To prove the lemma it suffices and is necessary to show that
$$
\text{length}_{\mathcal{O}_{X, \xi}}
(\mathcal{O}_{X, \xi}/a \mathcal{O}_{X, \xi}) =
\text{length}_{\mathcal{O}_{X, \xi}^h}
(\mathcal{O}_{X, \xi}^h/a \mathcal{O}_{X, \xi}^h)
$$
This follows immediately from
Algebra, Lemma \ref{algebra-lemma-pullback-module}
(and the fact that $\mathcal{O}_{X, \xi} \to \mathcal{O}_{X, \xi}^h$
is a flat local ring homomorphism of local Noetherian rings).
\end{proof}

\end{document}
